\section{Results and discussion}
\label{results}

\subsection{Classical techniques}

\subsubsection{Headline comparison and variability}

\textbf{Overall ranking.} Table~\ref{tab:classical-results} compares the classical pipelines over
five test partitions. Structured LR has the highest mean AP, $0.5928 \pm 0.1651$, followed by
text-only LR at $0.5762 \pm 0.1875$. Their mean F1 values are $0.4633$ and $0.5301$, so the
ranking by AP and the choice by thresholded F1 need not agree.

\begin{table}[tb]
\centering
\footnotesize
\setlength{\tabcolsep}{4pt}
\begin{tabular}{lccccr}
\toprule
\textbf{Classical model} & \textbf{Test AP} & \textbf{F1} & \textbf{Precision} & \textbf{Recall}
& \textbf{Seconds/split} \\
\midrule
\tabinput{includes/tab_classical_results}
\bottomrule
\end{tabular}
\caption{Classical-model test metrics over five grouped partitions (mean $\pm$ sample standard
deviation). F1, precision and recall use the validation-selected thresholds, except the majority
dummy's fixed threshold. These are saved elapsed times, not controlled end-to-end benchmarks.}
\label{tab:classical-results}
\end{table}

\textbf{Split variability.} Figure~\ref{fig:seeds} exposes the spread hidden by a mean score. All
four fine-tuned variants reach only 0.232 to 0.256 AP on seed 2, and several standard deviations in
Tables~\ref{tab:classical-results} and~\ref{tab:deep-results} exceed the between-family mean gaps.
One component built by the exact connector rule supplies 101 of the 171 fraud positives in seed-2
test data. This construction does not confirm that the component is one real-world fraud campaign.
Its concentration makes a nominal row-level score strongly dependent on one exact related-ad group,
so it should not be interpreted as 171 independent examples of fraud detection.

\begin{figure}[tb]
\centering
\includegraphics[width=0.94\linewidth]{test_ap_by_seed.pdf}
\caption{Test AP for each model on five grouped partitions. Points are individual runs, horizontal
ticks are means and the dashed line is mean fraud prevalence.}
\label{fig:seeds}
\end{figure}

\textbf{Validation/test behaviour.} Validation and test AP often disagree
(Figure~\ref{fig:validation-test}). This matters for transformer epoch selection as well as
classical tuning: a favourable validation trajectory need not transfer to the corresponding test
partition. Test results are reported for diagnosis and comparison, not used to select a more
favourable epoch or threshold.

\begin{figure}[tb]
\centering
\includegraphics[width=0.65\linewidth]{validation_vs_test_ap.pdf}
\caption{Validation and test AP for all 45 runs. The dashed line marks equal validation and test AP;
large departures motivate paired split-level analysis rather than validation-only conclusions.}
\label{fig:validation-test}
\end{figure}

\subsubsection{Classical structured-feature ablation}

\textbf{Effect direction.} Structured LR minus text-only LR has mean validation AP $+0.0333$ and
mean test AP $+0.0165$, with positive test differences on four of five splits. Seed 0 instead
changes by $-0.0473$. Mean test F1 changes by $-0.0668$ and improves on only one split. The paired
values in Table~\ref{tab:paired} and Figure~\ref{fig:ablation} therefore support a small,
split-dependent ranking effect, not a uniform improvement or a single best classical arm.

\begin{table}[tb]
\centering
\scriptsize
\setlength{\tabcolsep}{3pt}
\begin{tabular}{lrrrrrrc}
\toprule
\textbf{Metric difference} & \textbf{s0} & \textbf{s1} & \textbf{s2} & \textbf{s3}
& \textbf{s4} & \textbf{Mean} & \textbf{Positive} \\
\midrule
\tabinput{includes/tab_paired}
\bottomrule
\end{tabular}
\caption{Paired structured-LR-minus-text-only-LR differences. The final column counts strictly
positive changes among five partitions.}
\label{tab:paired}
\end{table}

\begin{figure}[tb]
\centering
\includegraphics[width=0.95\linewidth]{feature_ablation.pdf}
\caption{Per-partition effect of adding structured features to TF-IDF logistic regression. Positive
values favour structured features; legends report mean differences.}
\label{fig:ablation}
\end{figure}

\subsubsection{Confusion matrices}

\textbf{Seed-0 counts.} Figure~\ref{fig:confusion-classical} complements the classical AP/F1 results
with raw true-negative, false-positive, false-negative and true-positive counts. The majority dummy
classifies every ad as real; other panels expose how thresholds alter the error trade-off.

\begin{figure}[tb]
\centering
\includegraphics[width=0.98\linewidth]{confusion_classical.pdf}
\caption{Classical-model confusion matrices on the seed-0 test partition. Rows are actual labels,
columns are predicted labels, and cells contain TN, FP, FN and TP. Thresholds maximize validation F1,
except the majority dummy's fixed 0.5 threshold. One partition is diagnostic, not a five-split average.}
\label{fig:confusion-classical}
\end{figure}

\subsection{Deep learning techniques}

\subsubsection{Headline comparison}

\textbf{Overall ranking.} Table~\ref{tab:deep-results} compares text-only and structured-fusion
models under the two fine-tuning budgets. Full fused has the highest mean AP, $0.5047 \pm 0.1696$,
while full text reaches $0.4941 \pm 0.1645$. Their mean F1 values are $0.4540 \pm 0.2323$ and
$0.5247 \pm 0.1880$, again showing that AP and thresholded F1 need not select the same arm.

\begin{table}[tb]
\centering
\footnotesize
\setlength{\tabcolsep}{4pt}
\begin{tabular}{lccccr}
\toprule
\textbf{Deep-learning model} & \textbf{Test AP} & \textbf{F1} & \textbf{Precision} & \textbf{Recall}
& \textbf{Seconds/split} \\
\midrule
\tabinput{includes/tab_deep_results}
\bottomrule
\end{tabular}
\caption{Deep-learning test metrics over five grouped partitions (mean $\pm$ sample standard
deviation). F1, precision and recall use each run's validation-selected threshold. Seconds are saved
elapsed times, not a controlled hardware benchmark.}
\label{tab:deep-results}
\end{table}

\textbf{Comparison with structured TF-IDF LR.} Full text and full fused lose all five paired test-AP
comparisons to structured LR. Their respective mean AP deficits to that model are 0.0987 and 0.0881.
Full fine-tuning therefore does not outperform the strongest classical model in these paired runs.

\subsubsection{Adaptation budget and structured fusion}

\textbf{Full versus top two.} Table~\ref{tab:adaptation} reports paired full-minus-top-two changes.
Mean test AP increases by $+0.0245$ for text, with improvements on two splits, and by $+0.0474$ for
fusion, with improvements on three. The corresponding validation means are $+0.0745$ and $+0.0180$.
Updating the whole encoder therefore has small and inconsistent test benefits despite substantial
variation in individual splits.

\begin{table}[tb]
\centering
\scriptsize
\setlength{\tabcolsep}{4pt}
\begin{tabular}{lrrrcc}
\toprule
\textbf{Arm} & \textbf{Validation $\Delta$AP} & \textbf{Test $\Delta$AP}
& \textbf{Test $\Delta$F1} & \textbf{Test AP wins} & \textbf{Test F1 wins} \\
\midrule
\tabinput{includes/tab_adaptation}
\bottomrule
\end{tabular}
\caption{Paired full-minus-top-two differences over five partitions. Positive values and wins favour
full fine-tuning. Win columns report test partitions with strictly positive changes.}
\label{tab:adaptation}
\end{table}

\textbf{Late fusion.} Fusion effects also reverse between validation and test
(Table~\ref{tab:fusion}). For top-two tuning, mean validation AP rises by 0.0307 on all five splits,
but mean test AP falls by 0.0123 and rises on only two. For full tuning, fused-minus-text mean test
AP is $+0.0106$ with only two positive splits, while mean validation AP changes by $-0.0258$.
Full-model fusion also reduces mean test F1 by 0.0707. No stable fusion gain is established.

\begin{table}[tb]
\centering
\scriptsize
\setlength{\tabcolsep}{4pt}
\begin{tabular}{lrrrcc}
\toprule
\textbf{Strategy} & \textbf{Validation $\Delta$AP} & \textbf{Test $\Delta$AP}
& \textbf{Test $\Delta$F1} & \textbf{Test AP wins} & \textbf{Test F1 wins} \\
\midrule
\tabinput{includes/tab_fusion}
\bottomrule
\end{tabular}
\caption{Paired fused-minus-text differences for each adaptation strategy. Positive values and wins
favour structured late fusion.}
\label{tab:fusion}
\end{table}

\begin{figure}[tb]
\centering
\includegraphics[width=0.95\linewidth]{adaptation_ablation.pdf}
\caption{Paired full-minus-top-two AP differences for text-only and structured-fusion arms on
validation and test partitions. Positive values favour full encoder fine-tuning; legends report means.}
\label{fig:adaptation-ablation}
\end{figure}

\textbf{Saved run time.} Full text and full fused average approximately 155.0 and 155.7 seconds per
partition, compared with 106.4 and 107.2 seconds for the corresponding top-two runs. These records
describe the observed runs, not a controlled efficiency comparison; no common hardware benchmark
measures end-to-end development or serving cost across every family.

\subsubsection{Token coverage and operating points}

\textbf{Truncation audit.} On seed-0 training data, the 512-token cap covers 53.12\% of real and
88.72\% of fraud documents. The mean removed-token fraction over every document, assigning zero
removal to documents within the cap, is 13.10\% for real and 2.86\% for fraud ads. In the
approximate sequential field
audit among nonempty field entries, mean real-token survival is 68.05\% for requirements and 49.25\% for benefits, compared with
92.21\% and 92.58\% for fraud ads. Late-field loss is consequently class-dependent and descriptive,
not evidence that the removed text causes prediction errors.

\textbf{Thresholded profiles.} Figure~\ref{fig:metrics} separates precision, recall and F1 rather
than reducing every model to AP. Each threshold maximizes validation F1 independently, and wide F1
variation shows that this operating rule transfers unevenly across partitions.

\begin{figure}[tb]
\centering
\includegraphics[width=0.88\linewidth]{metric_profile.pdf}
\caption{Mean test precision, recall and F1 at validation-selected thresholds over five partitions.}
\label{fig:metrics}
\end{figure}

\textbf{Seed-0 curves.} Figure~\ref{fig:pr-curves} gives a score-level view for one already inspected
partition. It is a diagnostic illustration, not a basis for model selection; the means and paired
contrasts above carry the comparative evidence.

\begin{figure}[tb]
\centering
\includegraphics[width=0.86\linewidth]{pr_curves_seed0.pdf}
\caption{Seed-0 test precision-recall curves for the primary subset of six models: structured LR,
text-only LR, top-two text/fusion and full text/fusion. The dashed line is test fraud prevalence on
that partition.}
\label{fig:pr-curves}
\end{figure}

\subsubsection{Confusion matrices}

\textbf{Seed-0 counts.} Figure~\ref{fig:confusion-deep} shows the corresponding raw error counts for
all four fine-tuned variants. The panels compare false-alarm and missed-fraud trade-offs across
adaptation budgets and input types.

\begin{figure}[tb]
\centering
\includegraphics[width=0.98\linewidth]{confusion_deep.pdf}
\caption{Deep-learning confusion matrices on the seed-0 test partition. Rows are actual labels,
columns are predicted labels, and cells contain TN, FP, FN and TP. Every threshold is selected on
validation F1 only. These seed-0 counts are diagnostic and do not select a model or summarize all
five partitions.}
\label{fig:confusion-deep}
\end{figure}
